\section{Постановка задачи}
\label{sec:Chapter1}
\index{Chapter1}

Объектом исследования является подсистема SSA-оптимизаций (SSA --- Static Single Assignment форма представления инструкций) компилятора языка Go.

Предметом исследования являются методы автоматической векторизации ациклических участков кода на основе подхода Superword-Level Parallelism.

Целью работы является разработка и реализация прототипа механизма автоматической SLP-векторизации ациклических участков кода в компиляторе Go.

Для достижения поставленной цели в работе решаются следующие задачи:
\begin{itemize}
    \item исследовать существующие методы автоматической векторизации;
    \item проанализировать архитектуру оптимизационного конвейера компилятора Go;
    \item разработать алгоритм выявления векторизуемых групп операций;
    \item реализовать преобразования SSA-представления для генерации векторных операций;
    \item провести экспериментальную оценку корректности и эффективности разработанного решения.
\end{itemize}

Практическая значимость работы заключается в создании прототипа механизма автоматической векторизации для компилятора Go, 
демонстрирующего возможность интеграции SLP-преобразований в существующий компилятор. 
Полученные результаты могут быть использованы в качестве основы для дальнейшего развития оптимизирующих возможностей компилятора.